\section{公式、定理与例题}
\label{sec:00-math}

\subsection{行内公式与行间公式}

\begin{lstlisting}[language=TeX]
流水线的加速比约为 $S=\dfrac{n}{k+n-1}$，其中 $n$ 是指令条数，$k$ 是流水级数。

\[
  T_{\text{流水}} = \max_i t_i + (n-1)\cdot \max_i t_i
\]
\end{lstlisting}

效果：流水线的加速比约为 $S=\dfrac{n}{k+n-1}$，其中 $n$ 是指令条数，$k$ 是流水级数。

\[
  T_{\text{流水}} = \max_i t_i + (n-1)\cdot \max_i t_i
\]

行内公式前面不要留空行，\verb|$...$| 里的内容也不要换行。

\subsection{带编号、可引用的公式}

\begin{lstlisting}[language=TeX]
\begin{equation}
  \label{eq:00-speedup}
  S = \frac{T_{\text{串行}}}{T_{\text{流水}}}
\end{equation}
\end{lstlisting}

效果：

\begin{equation}
  \label{eq:00-speedup}
  S = \frac{T_{\text{串行}}}{T_{\text{流水}}}
\end{equation}

引用时写 \verb|\cref{eq:00-speedup}|，效果是\cref{eq:00-speedup}。需要换行对齐的
多行公式用 \texttt{align}：

\begin{lstlisting}[language=TeX]
\begin{align}
  S &= \frac{1}{1-p+\frac{p}{k}} \\
    &= \frac{k}{k+(1-p)k} \quad (\text{理想情况})
\end{align}
\end{lstlisting}

效果：

\begin{align}
  S &= \frac{1}{1-p+\frac{p}{k}} \\
    &= \frac{k}{k+(1-p)k} \quad (\text{理想情况})
\end{align}

\subsection{矩阵与分段函数}

\begin{lstlisting}[language=TeX]
\[
  A = \begin{pmatrix} a_{11} & a_{12} \\ a_{21} & a_{22} \end{pmatrix}
  \qquad
  |x| = \begin{cases}
    x,  & x \ge 0 \\
    -x, & x < 0
  \end{cases}
\]
\end{lstlisting}

\[
  A = \begin{pmatrix} a_{11} & a_{12} \\ a_{21} & a_{22} \end{pmatrix}
  \qquad
  |x| = \begin{cases}
    x,  & x \ge 0 \\
    -x, & x < 0
  \end{cases}
\]

\subsection{常用符号}

\begin{center}
\begin{tabular}{@{}llll@{}}
\toprule
写法 & 效果 & 写法 & 效果 \\
\midrule
\verb|\alpha, \beta, \theta| & $\alpha,\beta,\theta$ &
\verb|\sum_{i=1}^{n}| & $\sum_{i=1}^{n}$ \\
\verb|\int_0^T| & $\int_0^T$ &
\verb|\sqrt{x}| & $\sqrt{x}$ \\
\verb|\frac{a}{b}| & $\frac{a}{b}$ &
\verb|\cdot|、\verb|\times| & $\cdot$、$\times$ \\
\verb|\le|、\verb|\ge|、\verb|\ne| & $\le$、$\ge$、$\ne$ &
\verb|\in|、\verb|\subset| & $\in$、$\subset$ \\
\verb|\mathbb{R}| & $\mathbb{R}$ &
\verb|\mathcal{F}| & $\mathcal{F}$ \\
\bottomrule
\end{tabular}
\end{center}

\subsection{定义、定理、例题}

书里已经预置了三种环境，直接用：

\begin{lstlisting}[language=TeX]
\begin{definition}[指令集架构]
软件与硬件之间的契约，规定了可见的寄存器、指令语义、数据类型与异常行为；
它不规定这些行为由怎样的电路实现。
\end{definition}

\begin{theorem}[加速比上界]
在理想流水线下，加速比不超过流水级数 $k$。
\end{theorem}

\begin{example}
一条 5 级流水线，流水线填满后每周期完成 1 条指令。
\end{example}
\end{lstlisting}

效果：

\begin{definition}[指令集架构]
软件与硬件之间的契约，规定了可见的寄存器、指令语义、数据类型与异常行为；
它不规定这些行为由怎样的电路实现。
\end{definition}

\begin{theorem}[加速比上界]
在理想流水线下，加速比不超过流水级数 $k$。
\end{theorem}

\begin{example}
一条 5 级流水线，流水线填满后每周期完成 1 条指令。
\end{example}

证明用 \verb|\begin{proof}...\end{proof}|，结尾会自动画一个小方块。
